\chapter{Free Resolutions}
\label{ch:v23}

Free resolutions replace a module by an exact complex of free modules. They allow non-free modules to be studied using explicit matrices and are the computational input for Tor and Ext.

\section*{Learning goals}
After this chapter, the reader should be able to:
\begin{itemize}
\item construct free resolutions of elementary modules from generators and relations;
\item verify exactness of a proposed resolution at each stage;
\item compare two free resolutions through chain maps lifting the identity;
\item use resolutions to compute invariants that are independent of the chosen resolution;
\item truncate or extend resolutions while preserving exactness where justified;
\item distinguish a presentation from a full resolution by identifying higher syzygies;
\end{itemize}
\section*{Conceptual roadmap}
\[
\boxed{\text{presentation or universal property}\;\longrightarrow\;\text{exactness/localization}\;\longrightarrow\;\text{structural consequence}.}
\]

\section{Augmented complexes}
A resolution ends in a surjection $F_0\to M$ and is exact before and at $M$.

\section{Free resolutions}
All modules $F_i$ in a free resolution are free.

% BEGIN VOL05-EXPANSION V23-example-01
\begin{example}[Resolution of a cyclic group]\label{ex:v23-expand-01}
A free resolution of $\mathbb Z/n\mathbb Z$ begins $0\to\mathbb Z\xrightarrow{\cdot n}\mathbb Z\to\mathbb Z/n\mathbb Z\to0$. It is exact and both nonzero resolving modules are free.
\end{example}
% END VOL05-EXPANSION V23-example-01
\section{Construction by kernels}
Choose free generators for $M$, then free generators for the kernel, and repeat.

\section{Projective resolutions}
Projective resolutions generalize free resolutions and are often sufficient for derived functors.

\section{Comparison maps}
A module homomorphism lifts to a chain map between projective resolutions.

% BEGIN VOL05-EXPANSION V23-example-02
\begin{example}[Resolution of a quotient ring]\label{ex:v23-expand-02}
For a nonzerodivisor $f\in R$, the sequence $0\to R\xrightarrow{\cdot f}R\to R/(f)\to0$ is a free resolution of length one.
\end{example}
% END VOL05-EXPANSION V23-example-02
\section{Homotopy uniqueness}
Two such lifts are chain homotopic.

\section{Length}
Finite projective dimension means a resolution terminates after finitely many nonzero projectives.

% BEGIN VOL05-EXPANSION V23-example-03
\begin{example}[Computing after tensoring]\label{ex:v23-expand-03}
Applying a right-exact functor to a free resolution converts derived questions into homology calculations. Tensoring the resolution of $\mathbb Z/n$ with $\mathbb Z/m$ already reveals the kernel that becomes $\operatorname{Tor}_1$.
\end{example}
% END VOL05-EXPANSION V23-example-03
\section{Tensoring resolutions}
Tensoring a free resolution computes Tor homology.

\section{Applying Hom}
Applying Hom to a projective resolution computes Ext cohomology.

\section{Examples}
Cyclic quotients and regular sequences yield explicit short or structured resolutions.

\section{Core structural results}

\begin{theorem}[Existence of free resolutions]\label{thm:v23-01}
Every module admits a free resolution.
\begin{proof}
Choose a free module surjecting onto the module. Take its kernel, choose a free module surjecting onto that kernel, and iterate indefinitely.
\end{proof}
\end{theorem}

\begin{theorem}[Comparison theorem]\label{thm:v23-02}
A homomorphism $M\to N$ lifts to a chain map from any projective resolution of $M$ to any resolution of $N$ ending in a surjection.
\begin{proof}
Lift degree zero using projectivity. Inductively, the image of the previous differential lands in the next kernel by the chain condition, and projectivity lifts through the corresponding surjection.
\end{proof}
\end{theorem}

\begin{theorem}[Homotopy uniqueness of comparison maps]\label{thm:v23-03}
Two chain maps between projective resolutions inducing the same module map are chain homotopic.
\begin{proof}
Construct the homotopy inductively. At each degree, the difference of the two maps has image in a kernel already identified with the image of the next differential, and projectivity provides a lift.
\end{proof}
\end{theorem}

\section{Worked examples}

\begin{example}[Augmented complexes diagnostic]\label{ex:v23-worked-01}
Give a rigorous worked analysis of Augmented complexes. State the ring, module, finiteness, localization, or homological hypotheses and derive the central conclusion. A resolution ends in a surjection $F_0\to M$ and is exact before and at $M$. The decisive step is to use the universal property, exactness criterion, localization test, or finite-generation mechanism appropriate to the algebraic structure.
\end{example}

\begin{example}[Free resolutions diagnostic]\label{ex:v23-worked-02}
Give a rigorous worked analysis of Free resolutions. State the ring, module, finiteness, localization, or homological hypotheses and derive the central conclusion. All modules $F_i$ in a free resolution are free. The decisive step is to use the universal property, exactness criterion, localization test, or finite-generation mechanism appropriate to the algebraic structure.
\end{example}

\begin{example}[Construction by kernels diagnostic]\label{ex:v23-worked-03}
Give a rigorous worked analysis of Construction by kernels. State the ring, module, finiteness, localization, or homological hypotheses and derive the central conclusion. Choose free generators for $M$, then free generators for the kernel, and repeat. The decisive step is to use the universal property, exactness criterion, localization test, or finite-generation mechanism appropriate to the algebraic structure.
\end{example}

\begin{example}[Projective resolutions diagnostic]\label{ex:v23-worked-04}
Give a rigorous worked analysis of Projective resolutions. State the ring, module, finiteness, localization, or homological hypotheses and derive the central conclusion. Projective resolutions generalize free resolutions and are often sufficient for derived functors. The decisive step is to use the universal property, exactness criterion, localization test, or finite-generation mechanism appropriate to the algebraic structure.
\end{example}

\section{Solved dossiers}
Each dossier is a canonical solved problem. The provenance ledger records which retained corpus rules guided its topic and which dossiers were added freshly to complete the chapter.

\begin{problem}[Augmented complexes diagnostic]\label{prob:v23-dossier-01}
Give a rigorous worked analysis of Augmented complexes. State the ring, module, finiteness, localization, or homological hypotheses and derive the central conclusion.
\end{problem}
\begin{solution}
A resolution ends in a surjection $F_0\to M$ and is exact before and at $M$. The decisive step is to use the universal property, exactness criterion, localization test, or finite-generation mechanism appropriate to the algebraic structure.
\end{solution}

\begin{problem}[Free resolutions diagnostic]\label{prob:v23-dossier-02}
Give a rigorous worked analysis of Free resolutions. State the ring, module, finiteness, localization, or homological hypotheses and derive the central conclusion.
\end{problem}
\begin{solution}
All modules $F_i$ in a free resolution are free. The decisive step is to use the universal property, exactness criterion, localization test, or finite-generation mechanism appropriate to the algebraic structure.
\end{solution}

\begin{problem}[Construction by kernels diagnostic]\label{prob:v23-dossier-03}
Give a rigorous worked analysis of Construction by kernels. State the ring, module, finiteness, localization, or homological hypotheses and derive the central conclusion.
\end{problem}
\begin{solution}
Choose free generators for $M$, then free generators for the kernel, and repeat. The decisive step is to use the universal property, exactness criterion, localization test, or finite-generation mechanism appropriate to the algebraic structure.
\end{solution}

\begin{problem}[Projective resolutions diagnostic]\label{prob:v23-dossier-04}
Give a rigorous worked analysis of Projective resolutions. State the ring, module, finiteness, localization, or homological hypotheses and derive the central conclusion.
\end{problem}
\begin{solution}
Projective resolutions generalize free resolutions and are often sufficient for derived functors. The decisive step is to use the universal property, exactness criterion, localization test, or finite-generation mechanism appropriate to the algebraic structure.
\end{solution}

\begin{problem}[Comparison maps diagnostic]\label{prob:v23-dossier-05}
Give a rigorous worked analysis of Comparison maps. State the ring, module, finiteness, localization, or homological hypotheses and derive the central conclusion.
\end{problem}
\begin{solution}
A module homomorphism lifts to a chain map between projective resolutions. The decisive step is to use the universal property, exactness criterion, localization test, or finite-generation mechanism appropriate to the algebraic structure.
\end{solution}

\begin{problem}[Homotopy uniqueness diagnostic]\label{prob:v23-dossier-06}
Give a rigorous worked analysis of Homotopy uniqueness. State the ring, module, finiteness, localization, or homological hypotheses and derive the central conclusion.
\end{problem}
\begin{solution}
Two such lifts are chain homotopic. The decisive step is to use the universal property, exactness criterion, localization test, or finite-generation mechanism appropriate to the algebraic structure.
\end{solution}

\begin{problem}[Length diagnostic]\label{prob:v23-dossier-07}
Give a rigorous worked analysis of Length. State the ring, module, finiteness, localization, or homological hypotheses and derive the central conclusion.
\end{problem}
\begin{solution}
Finite projective dimension means a resolution terminates after finitely many nonzero projectives. The decisive step is to use the universal property, exactness criterion, localization test, or finite-generation mechanism appropriate to the algebraic structure.
\end{solution}

\begin{problem}[Tensoring resolutions diagnostic]\label{prob:v23-dossier-08}
Give a rigorous worked analysis of Tensoring resolutions. State the ring, module, finiteness, localization, or homological hypotheses and derive the central conclusion.
\end{problem}
\begin{solution}
Tensoring a free resolution computes Tor homology. The decisive step is to use the universal property, exactness criterion, localization test, or finite-generation mechanism appropriate to the algebraic structure.
\end{solution}

\begin{problem}[Existence of free resolutions proof dossier]\label{prob:v23-dossier-09}
Prove the following structural statement and identify the hypothesis that makes the conclusion work: Every module admits a free resolution.
\end{problem}
\begin{solution}
Choose a free module surjecting onto the module. Take its kernel, choose a free module surjecting onto that kernel, and iterate indefinitely. This proof isolates the reusable algebraic mechanism behind the result.
\end{solution}

\begin{problem}[Comparison theorem proof dossier]\label{prob:v23-dossier-10}
Prove the following structural statement and identify the hypothesis that makes the conclusion work: A homomorphism $M\to N$ lifts to a chain map from any projective resolution of $M$ to any resolution of $N$ ending in a surjection.
\end{problem}
\begin{solution}
Lift degree zero using projectivity. Inductively, the image of the previous differential lands in the next kernel by the chain condition, and projectivity lifts through the corresponding surjection. This proof isolates the reusable algebraic mechanism behind the result.
\end{solution}

\begin{problem}[Homotopy uniqueness of comparison maps proof dossier]\label{prob:v23-dossier-11}
Prove the following structural statement and identify the hypothesis that makes the conclusion work: Two chain maps between projective resolutions inducing the same module map are chain homotopic.
\end{problem}
\begin{solution}
Construct the homotopy inductively. At each degree, the difference of the two maps has image in a kernel already identified with the image of the next differential, and projectivity provides a lift. This proof isolates the reusable algebraic mechanism behind the result.
\end{solution}

\begin{problem}[Chapter synthesis]\label{prob:v23-dossier-12}
Connect the definitions and principal structural theorems of this chapter into one reusable commutative-algebra strategy.
\end{problem}
\begin{solution}
Free resolutions replace a module by an exact complex of free modules. They allow non-free modules to be studied using explicit matrices and are the computational input for Tor and Ext. The recurring strategy is to encode the algebraic object by kernels, quotients, localization, finite presentation, or a resolution, apply the corresponding universal or exactness principle, and then recover the desired global or local conclusion.
\end{solution}

\section{Exercises with complete solutions}

\begin{exercise}\label{exr:v23-01}
State and apply the central principle behind Augmented complexes. Give one concrete algebraic consequence.
\end{exercise}
\begin{hint}
Start from a surjection \(F_0\to M\) from a free module on chosen generators; the first syzygy is its kernel.
\end{hint}
\begin{solution}
A resolution ends in a surjection $F_0\to M$ and is exact before and at $M$. The same principle can then be reused in later localization, support, base-change, resolution, Tor, or Ext arguments.
\end{solution}

\begin{exercise}\label{exr:v23-02}
State and apply the central principle behind Free resolutions. Give one concrete algebraic consequence.
\end{exercise}
\begin{hint}
Choose a free module \(F_1\) surjecting onto that kernel, then repeat recursively to build a free resolution.
\end{hint}
\begin{solution}
All modules $F_i$ in a free resolution are free. The same principle can then be reused in later localization, support, base-change, resolution, Tor, or Ext arguments.
\end{solution}

\begin{exercise}\label{exr:v23-03}
State and apply the central principle behind Construction by kernels. Give one concrete algebraic consequence.
\end{exercise}
\begin{hint}
At each degree, exactness requires image of the new differential to equal the previous kernel.
\end{hint}
\begin{solution}
Choose free generators for $M$, then free generators for the kernel, and repeat. The same principle can then be reused in later localization, support, base-change, resolution, Tor, or Ext arguments.
\end{solution}

\begin{exercise}\label{exr:v23-04}
State and apply the central principle behind Projective resolutions. Give one concrete algebraic consequence.
\end{exercise}
\begin{hint}
A presentation stops after \(F_1\to F_0\to M\); a resolution continues by resolving the relation modules.
\end{hint}
\begin{solution}
Projective resolutions generalize free resolutions and are often sufficient for derived functors. The same principle can then be reused in later localization, support, base-change, resolution, Tor, or Ext arguments.
\end{solution}

\begin{exercise}\label{exr:v23-05}
State and apply the central principle behind Comparison maps. Give one concrete algebraic consequence.
\end{exercise}
\begin{hint}
Comparison maps between resolutions are built inductively using projectivity of the free modules.
\end{hint}
\begin{solution}
A module homomorphism lifts to a chain map between projective resolutions. The same principle can then be reused in later localization, support, base-change, resolution, Tor, or Ext arguments.
\end{solution}

\begin{exercise}\label{exr:v23-06}
State and apply the central principle behind Homotopy uniqueness. Give one concrete algebraic consequence.
\end{exercise}
\begin{hint}
Different resolutions may look different, so only use invariants proved independent up to chain homotopy.
\end{hint}
\begin{solution}
Two such lifts are chain homotopic. The same principle can then be reused in later localization, support, base-change, resolution, Tor, or Ext arguments.
\end{solution}

\begin{exercise}\label{exr:v23-07}
State and apply the central principle behind Length. Give one concrete algebraic consequence.
\end{exercise}
\begin{hint}
Truncating a resolution changes exactness at the cut unless the remaining kernel is projective/free as required.
\end{hint}
\begin{solution}
Finite projective dimension means a resolution terminates after finitely many nonzero projectives. The same principle can then be reused in later localization, support, base-change, resolution, Tor, or Ext arguments.
\end{solution}

\begin{exercise}\label{exr:v23-08}
State and apply the central principle behind Tensoring resolutions. Give one concrete algebraic consequence.
\end{exercise}
\begin{hint}
For explicit modules, derive the first differential from a presentation matrix before inventing higher maps.
\end{hint}
\begin{solution}
Tensoring a free resolution computes Tor homology. The same principle can then be reused in later localization, support, base-change, resolution, Tor, or Ext arguments.
\end{solution}

\section*{Chapter summary}
The chapter now contributes a canonical solved layer to the progression from rings and modules through localization, Noetherian methods, integral dependence, and derived functors.
% BEGIN VOL05-EXPANSION V23-exercises-01
\section{Graded supplementary exercises}
These exercises extend the protected chapter with a balanced set of computations, proofs, hypothesis tests, applications, and challenges.

\subsection*{Standard computations}

\begin{exercise}[Cyclic resolution]\label{exr:v23-expand-01}
Write a free resolution of $\mathbb Z/5$ of length one.
\end{exercise}
\begin{hint}
Use free resolutions and exact augmentations.
\end{hint}
\begin{solution}
$0\to\mathbb Z\xrightarrow{\cdot5}\mathbb Z\to\mathbb Z/5\to0$.
\end{solution}

\begin{exercise}[Quotient resolution]\label{exr:v23-expand-02}
Resolve $R/(f)$ when $f$ is a nonzerodivisor.
\end{exercise}
\begin{hint}
Use free resolutions and exact augmentations.
\end{hint}
\begin{solution}
$0\to R\xrightarrow{\cdot f}R\to R/(f)\to0$.
\end{solution}

\begin{exercise}[Projective dimension]\label{exr:v23-expand-03}
What upper bound follows from a length-one free resolution?
\end{exercise}
\begin{hint}
Use free resolutions and exact augmentations.
\end{hint}
\begin{solution}
Projective dimension at most one.
\end{solution}

\begin{exercise}[Free module]\label{exr:v23-expand-04}
What resolution does a free module need?
\end{exercise}
\begin{hint}
Use free resolutions and exact augmentations.
\end{hint}
\begin{solution}
Length zero.
\end{solution}

\begin{exercise}[Augmentation]\label{exr:v23-expand-05}
What is the last map in a resolution called?
\end{exercise}
\begin{hint}
Use free resolutions and exact augmentations.
\end{hint}
\begin{solution}
The augmentation onto the resolved module.
\end{solution}

\subsection*{Proofs}

\begin{exercise}[Existence of free resolutions proof]\label{exr:v23-expand-06}
Prove: Every module admits a free resolution.
\end{exercise}
\begin{hint}
Start from the theorem hypotheses and identify the decisive algebraic construction.
\end{hint}
\begin{solution}
Choose a free module surjecting onto the module. Take its kernel, choose a free module surjecting onto that kernel, and iterate indefinitely.
\end{solution}

\begin{exercise}[Comparison theorem proof]\label{exr:v23-expand-07}
Prove: A homomorphism $M\to N$ lifts to a chain map from any projective resolution of $M$ to any resolution of $N$ ending in a surjection.
\end{exercise}
\begin{hint}
Start from the theorem hypotheses and identify the decisive algebraic construction.
\end{hint}
\begin{solution}
Lift degree zero using projectivity. Inductively, the image of the previous differential lands in the next kernel by the chain condition, and projectivity lifts through the corresponding surjection.
\end{solution}

\begin{exercise}[Homotopy uniqueness of comparison maps proof]\label{exr:v23-expand-08}
Prove: Two chain maps between projective resolutions inducing the same module map are chain homotopic.
\end{exercise}
\begin{hint}
Start from the theorem hypotheses and identify the decisive algebraic construction.
\end{hint}
\begin{solution}
Construct the homotopy inductively. At each degree, the difference of the two maps has image in a kernel already identified with the image of the next differential, and projectivity provides a lift.
\end{solution}

\begin{exercise}[Structural synthesis]\label{exr:v23-expand-09}
Prove a corollary that genuinely uses both Existence of free resolutions and Comparison theorem.
\end{exercise}
\begin{hint}
Apply the first theorem to obtain its structural description, then verify the second theorem's hypotheses.
\end{hint}
\begin{solution}
A correct proof explicitly invokes both results, verifies the common hypotheses, and derives a new consequence rather than merely restating either theorem.
\end{solution}

\subsection*{Counterexamples and hypothesis tests}

\begin{exercise}[Unique resolution]\label{exr:v23-expand-10}
Test the claim: A module has a unique free resolution.
\end{exercise}
\begin{hint}
Use the smallest concrete ring, module, or resolution that can decide the claim.
\end{hint}
\begin{solution}
False; resolutions are highly nonunique.
\end{solution}

\begin{exercise}[Finite resolution]\label{exr:v23-expand-11}
Test the claim: Every module has a finite free resolution.
\end{exercise}
\begin{hint}
Use the smallest concrete ring, module, or resolution that can decide the claim.
\end{hint}
\begin{solution}
False in general.
\end{solution}

\begin{exercise}[Resolution exactness]\label{exr:v23-expand-12}
Test the claim: A resolution may have nonzero homology in positive degrees.
\end{exercise}
\begin{hint}
Use the smallest concrete ring, module, or resolution that can decide the claim.
\end{hint}
\begin{solution}
False; it is exact there.
\end{solution}

\subsection*{Applications and investigations}

\begin{exercise}[Derived computation]\label{exr:v23-expand-13}
Why resolve by free or projective modules?
\end{exercise}
\begin{hint}
Translate the question into free resolutions and exact augmentations.
\end{hint}
\begin{solution}
Such modules make Hom and tensor computations tractable and define derived functors.
\end{solution}

\begin{exercise}[Complexity]\label{exr:v23-expand-14}
What does minimal resolution length measure?
\end{exercise}
\begin{hint}
Translate the question into free resolutions and exact augmentations.
\end{hint}
\begin{solution}
It contributes to projective dimension and homological complexity.
\end{solution}

\subsection*{Challenge problems}

\begin{exercise}[Local-global synthesis]\label{exr:v23-expand-15}
Connect 'Augmented complexes' with 'Examples' in one rigorous argument.
\end{exercise}
\begin{hint}
State the relevant ring map, module map, or complex before applying the structural theorem.
\end{hint}
\begin{solution}
The opening idea is: A resolution ends in a surjection $F_0\to M$ and is exact before and at $M$. The later viewpoint is: Cyclic quotients and regular sequences yield explicit short or structured resolutions. A complete solution identifies the structural theorem that links these descriptions and checks its hypotheses.
\end{solution}

\begin{exercise}[Hypothesis audit]\label{exr:v23-expand-16}
Choose one theorem from the chapter, remove one essential hypothesis, and exhibit or explain a concrete failure.
\end{exercise}
\begin{hint}
Use one of the chapter's explicit computations or counterexamples as the test case.
\end{hint}
\begin{solution}
The solution must name the removed hypothesis, show exactly which proof step breaks, and give a concrete algebraic object where the claimed conclusion fails.
\end{solution}

% END VOL05-EXPANSION V23-exercises-01
